% 对应题号 M1-66，复用已核对的本题几何图。
% 核心结论：冲量等于动量跃变量，也等于集中力分布的权重。
% 图契约：物理竞赛讲义；印刷宽16cm，最小字号8.5pt；几何示意与公式分区。
% 布局：[几何/曲线 7cm] [定义/结果 7cm]；实线主量，虚线辅助；无装饰数据。
\documentclass[tikz,border=6pt]{standalone}
\usepackage{fontspec,xeCJK,amsmath,amssymb,graphicx}
\IfFontExistsTF{Times New Roman}{\setmainfont{Times New Roman}}{\setmainfont{TeX Gyre Termes}}
\IfFontExistsTF{Heiti SC}{\setCJKmainfont[AutoFakeBold=1.5]{Heiti SC}}{\setCJKmainfont[AutoFakeBold=1.5]{FandolHei-Regular}}
\usetikzlibrary{arrows.meta,bending,calc,positioning,patterns}
\definecolor{ink}{HTML}{243745}
\definecolor{blue}{HTML}{275DAB}
\definecolor{green}{HTML}{177C63}
\definecolor{red}{HTML}{BD4935}
\definecolor{orange}{HTML}{AA6C16}
\definecolor{purple}{HTML}{715393}
\tikzset{
 every node/.style={font=\fontsize{9.5}{13}\selectfont,text=ink},
 axis/.style={-{Stealth[length=4pt]},line width=.65pt,draw=ink!65},
 vec/.style={-{Stealth[length=5pt]},line width=1pt,draw=blue},
 short/.style={-{Stealth[length=3pt]},line width=.8pt},
 guide/.style={dashed,draw=ink!40,line width=.5pt},
 card/.style={rounded corners=2pt,draw=blue!60,fill=blue!4,
   align=center,text width=6.8cm,inner xsep=9pt,inner ysep=7pt},
 title/.style={font=\fontsize{11}{14}\selectfont\bfseries,text width=15cm,align=center},
 label/.style={font=\fontsize{8.5}{11}\selectfont,align=center}}
\begin{document}
\begin{tikzpicture}[x=1cm,y=1cm]
\path[use as bounding box] (-8,-3.35) rectangle (8,3.35);
\node[title] at (0,2.95) {例：瞬时冲量对应动量的跃变};
\begin{scope}[shift={(-4,0)}]

\draw[axis] (-2.5,-1.4)--(2.5,-1.4) node[right] {$t$};\draw[axis] (-2.2,-1.7)--(-2.2,2) node[above] {$p$};
\draw[blue,line width=1pt] (-2,-.7)--(0,-.7);\draw[blue,line width=1pt] (0,1.2)--(2.1,1.2);\draw[guide] (0,-.7)--(0,1.2);
\draw[<->,red] (.4,-.7)--(.4,1.2) node[midway,right] {$I$};\node[below] at (0,-1.4) {$t_0$};
\node[label] at (0,-2.45) {纵轴为某一固定方向上的动量分量};

\end{scope}\node[card,draw=blue!65,fill=blue!4] at (3.9,1.65) {$p(t)=p_-+I H(t-t_0)$};
\node[card,draw=green!65,fill=green!4] at (3.9,0) {$\displaystyle\frac{dp}{dt}=I\delta(t-t_0)$};
\node[card,draw=orange!65,fill=orange!4] at (3.9,-1.65) {$p_+-p_-=\int F\,dt=I$\\δ 的权重等于总冲量。};
\end{tikzpicture}\end{document}
